\documentclass[10pt,a4paper]{amsart}
\usepackage[margin=19mm]{geometry}
\usepackage{amsmath,amssymb,mathtools}
\usepackage[scr=rsfs,cal=euler]{mathalfa}
\usepackage{xfrac}
\usepackage[unicode,hidelinks,bookmarks=false]{hyperref}
\newcommand{\ie}{i.e.\ }
\newcommand{\cf}{cf.\ }
\newcommand{\resp}{respectively\ }
\newcommand{\Subsection}{\subsection}
\providecommand{\nobreakdash}{\nobreak\hbox{-}}
\setlength{\emergencystretch}{2em}
\multlinegap0pt
\usepackage{amssymb}
\usepackage{tikz}
\usepackage{float}
\usepackage{algorithm}\usepackage{algpseudocode}
\newtheorem{theorem}{Theorem}
\newcommand{\LT}{\mathrm{L}}
\newcommand{\LaT}{\mathrm{La}}
\newcommand{\floor}[1]{\flr{#1}}
\newcommand{\flr}[1]{\left\lfloor #1 \right\rfloor}
\title{\textbf{Erd\H{o}s-Gy\'{a}rf\'{a}s problem for partially ordered sets}}
\author{Gyula O.\, H. Katona \and Yaping Mao}
\date{Source: arXiv:2604.10229v2 (CC BY 4.0)}
\begin{document}
\maketitle

\noindent\emph{Source and licence.} \textbf{Erd\H{o}s-Gy\'{a}rf\'{a}s problem for partially ordered sets}. Version arXiv:2604.10229v2 (CC BY 4.0), distributed by arXiv under the Creative Commons Attribution 4.0 International licence, http://creativecommons.org/licenses/by/4.0/, as recorded in the arXiv licence metadata for this exact version and shown on the abstract page https://arxiv.org/abs/2604.10229v2. Full licence notice: \texttt{licenses/arxiv-2604.10229-CC-BY-4.0.txt}.\par\smallskip

\noindent\textit{Editorial scope.} Counted unit: the article's theorem th-Lat (source0.tex lines 776--791). The article's complete original proof of it is delivered with it (source0.tex lines 793--842, one \texttt{proof} environment of 50 lines). No mathematics is written, repaired or re-derived: the statement and the proof are the article's own lines, in the article's own notation, and no retained line is substituted or re-flowed.  No further source-local context is needed: every cross-reference of every retained line resolves inside the two delivered spans.  Nothing was omitted from any retained span, so the package carries no omission record.  No retained line contains a citation, so the wrapper carries no bibliography and no bibliography entry.  No retained line contains a figure environment, an image-inclusion command or drawing code, so no image is delivered and the package declares no asset dependency.  Editorial formatting only, in addition: an AMS \texttt{amsart} preamble that restates the article's own declarations and macros used by the retained lines, namely the environments \texttt{theorem} on separate counters, together with the standard packages those lines need.  The retained environments print in this document as Theorem 1 in order of appearance, whereas the article prints the same items with the numbering of its own full document.  The complete native source of the article as distributed in the arXiv e-print for this exact version is bundled unchanged as \texttt{sources/198252/source0.tex}.
\begin{theorem}\label{th-Lat}
Let $m\ge2$, $1\le t\le m$, and $n\ge m+1$. Then
\[
h_{n,m}(t)
\le
\LaT_t^\sharp(n,B_m)
\le
h_n(t)-\floor{\frac{n}{m+1}}\binom{m}{t-1},
\]
and
\[
\LT_t^\sharp(n,B_m)
\le
h_n(t)-\floor{\frac{n}{m+1}}.
\]
\end{theorem}

\begin{proof}
Let
\[
\mathcal M_{n,m}
=
\bigcup_{k=r}^{r+m-1}\mathcal S_k(B_n),
\qquad
r=\floor{\frac{n-m}{2}},
\]
where \(\mathcal S_k(B_n)=\{X\subseteq[n]:|X|=k\}\). Since \(\mathcal M_{n,m}\) has height \(m\), it contains no induced copy of \(B_m\), whose height is \(m+1\). Hence
$\LaT_t^\sharp(n,B_m)\ge h_{n,m}(t)$.

Set $s=\floor{\frac{n}{m+1}}$.
Partition the first \((m+1)s\) coordinates into disjoint blocks
\[
J_k=\{(k-1)(m+1)+1,\dots,k(m+1)\},
\qquad 1\le k\le s.
\]
Fix a distinguished element \(d_k\in J_k\) and write \(I_k=J_k\setminus\{d_k\}\). Define
\[
Q_k=\bigl\{\{d_k\}\cup A:A\subseteq I_k\bigr\}.
\]
Each \(Q_k\) is an induced copy of \(B_m\), and the families \(Q_1,\dots,Q_s\) are pairwise disjoint.

To prove the upper bound on \(\LaT_t^\sharp(n,B_m)\), let \(\mathcal F\subseteq B_n\) be strongly \(B_m\)-free. Then \(\mathcal F\cap Q_k\neq Q_k\) for every \(k\), so choose \(x_k\in Q_k\setminus \mathcal F\). Write \(x_k=\{d_k\}\cup A_k\) with \(|A_k|=r_k\). Fix an arbitrary linear ordering of \(I_k\). For every choice of subsets
\[
D\subseteq A_k,\qquad U\subseteq I_k\setminus A_k,\qquad |D|+|U|=t-1,
\]
we obtain a \(t\)-chain in \(Q_k\) through \(x_k\) by deleting the elements of \(D\) from \(x_k\) one by one in the chosen order and then adding the elements of \(U\) one by one in the chosen order. Distinct pairs \((D,U)\) yield distinct \(t\)-chains. Hence \(x_k\) belongs to at least
\[
\sum_{i=0}^{t-1}\binom{r_k}{i}\binom{m-r_k}{t-1-i}
=
\binom{m}{t-1}
\]
distinct \(t\)-chains of \(Q_k\), by Vandermonde's identity.
All these \(t\)-chains are absent from \(A_t(\mathcal F)\), and because the \(Q_k\) are pairwise disjoint, the omitted \(t\)-chains coming from different \(k\) are distinct. Therefore
\[
|A_t(\mathcal F)|\le h_n(t)-s\binom{m}{t-1}.
\]
Taking the maximum over all strongly \(B_m\)-free families \(\mathcal F\) yields
\[
\LaT_t^\sharp(n,B_m)\le h_n(t)-\floor{\frac{n}{m+1}}\binom{m}{t-1}.
\]

To prove the upper bound on \(\LT_t^\sharp(n,B_m)\), let \(\mathcal C_t\subseteq \mathcal T_n(t)\) be strongly \((t,B_m)\)-free. Then \(A_t(Q_k)\nsubseteq \mathcal C_t\) for every \(k\), so choose \(T_k\in A_t(Q_k)\setminus \mathcal C_t\). Since the copies \(Q_1,\dots,Q_s\) are pairwise disjoint, the sets \(A_t(Q_k)\) are pairwise disjoint, and thus the omitted chains \(T_1,\dots,T_s\) are distinct. Hence
\[
|\mathcal C_t|\le h_n(t)-s=h_n(t)-\floor{\frac{n}{m+1}}.
\]
Taking the maximum over all strongly \((t,B_m)\)-free sets \(\mathcal C_t\) gives the desired bound.
\end{proof}

\end{document}
